\chapter{Wo die Zeit hingeht}\label{kap:zeit_hin}

\begin{merksatz}
\boxtitel{Die Aufteilung einer Anfrage}
\begin{center}
{\small
\renewcommand{\arraystretch}{1.5}
\begin{tabular}{@{}P{3.0cm}P{2.0cm}P{2.6cm}@{}}
\toprule
\textbf{Abschnitt} & \textbf{Anteil} & \\
\midrule
\textbf{Datenbank\-abfragen} & $\approx\mathbf{48\,\%}$ & \textcolor{akzent}{\textbf{davon $n{+}1$: $32\,\%$}}\\
Netz zwischen Diensten & $\approx18\,\%$ & \\
Serialisierung & $\approx12\,\%$ & \textcolor{akzent}{oft übersehen}\\
Eigene Rechnung & $\approx11\,\%$ & \\
Warten auf Sperren & $\approx7\,\%$ & \\
Protokollierung & $\approx4\,\%$ & \textbf{Band 22, Kap.~12}\\
\bottomrule
\end{tabular}}
\end{center}

\begin{center}
\begin{minipage}{0.88\textwidth}
\kursiv{Elf Prozent der Zeit vergehen dort, wo die meisten Optimierungsversuche ansetzen. Fast die Hälfte in der Datenbank, und zwei Drittel davon durch ein einziges Muster --- dasselbe $n{+}1$-Problem wie in Band 20, Kapitel~12 und Band 23, Kapitel~\ref{kap:naht}.}
\end{minipage}
\end{center}

\textbf{Die Serialisierung, die niemand rechnet.}

\begin{center}
\renewcommand{\arraystretch}{1.45}
\begin{tabular}{@{}P{2.6cm}P{2.2cm}P{2.4cm}P{2.0cm}@{}}
\toprule
\textbf{Vorgang} & \textbf{je $1000$} & \textbf{Bei $5$ Nähten} & \\
\midrule
JSON erzeugen & $4{,}8$ ms & $\mathbf{24}$ ms & \\
JSON lesen & $6{,}1$ ms & $\mathbf{31}$ ms & \\
Binär erzeugen & $\mathbf{0{,}2}$ ms & $1$ ms & $24\times$\\
Binär lesen & $0{,}3$ ms & $1{,}5$ ms & \\
\bottomrule
\end{tabular}
\end{center}

\begin{center}
\kursiv{Fünfundfünfzig Millisekunden für das Umwandeln derselben tausend Befunde, wenn sie durch fünf Nähte gehen. Das ist mehr, als die Fachlogik braucht --- und es steht in keinem Profil eines einzelnen Dienstes.}
\end{center}

\begin{technik}
\boxtitel{Die Ablaufverfolgung über Nähte}
\begin{center}
\renewcommand{\arraystretch}{1.5}
\begin{tabular}{@{}P{2.6cm}P{6.0cm}@{}}
\toprule
\textbf{Das Mittel} & Eine Kennung wandert mit der Anfrage durch alle Bestandteile\\
\textbf{Was man sieht} & \textbf{Wo die Zeit vergeht --- über Dienstgrenzen hinweg}\\
\textbf{Der Aufwand} & Kennung durchreichen, Abschnitte melden\\
\bottomrule
\end{tabular}
\end{center}

\begin{center}
\begin{minipage}{0.86\textwidth}
\kursiv{Ohne dieses Mittel misst jeder Dienst sich selbst, und niemand sieht das Ganze. Mit ihm bekommt man ein Bild einer einzelnen Anfrage über alle Nähte --- und die Frage \enquote{warum hat das zwei Sekunden gedauert} wird beantwortbar.}
\end{minipage}
\end{center}

\begin{center}
\kursiv{Der Aufwand ist gering und die Disziplin hoch: Die Kennung muss durch jede Naht weitergereicht werden, auch durch Warteschlangen und asynchrone Verarbeitung. Wo sie verlorengeht, endet die Verfolgung.}
\end{center}
\end{technik}
\end{merksatz}

\begin{merksatz}
\boxtitel{Das $n{+}1$-Problem über Nähte}
\begin{center}
\begin{minipage}{0.9\textwidth}
\code{// Auswertung fragt die Speicherung}\\
\code{faelle = hole\_faelle(patient)\ \ \ \ \ \ \ // 1 Aufruf}\\
\code{for fall in faelle:}\\
\code{\ \ \ \ befunde = hole\_befunde(fall.id)\ // n Aufrufe ueber die Naht}
\end{minipage}
\end{center}

\begin{center}
\renewcommand{\arraystretch}{1.45}
\begin{tabular}{@{}P{2.2cm}P{2.2cm}P{2.2cm}P{2.2cm}@{}}
\toprule
\textbf{Fälle} & \textbf{Im Prozess} & \textbf{Über Naht} & \textbf{Faktor}\\
\midrule
$12$ & $52$ ms & $\mathbf{190}$ ms & $3{,}7$\\
$80$ & $324$ ms & $1{,}2$ s & \\
$340$ & $1{,}4$ s & \textcolor{akzent}{$\mathbf{5{,}1}$ \textbf{s}} & \textcolor{akzent}{\textbf{Langzeit\-patient}}\\
\bottomrule
\end{tabular}
\end{center}

\begin{center}
\begin{minipage}{0.88\textwidth}
\kursiv{Dasselbe Muster wie in Band 20, mit vierfachem Faktor --- weil jeder Aufruf jetzt Netz, Serialisierung und Warteschlange durchläuft. Fünf Sekunden für den Aufbau einer Patientenansicht, und das bei dem Patienten, bei dem es am meisten zählt.}
\end{minipage}
\end{center}

\textbf{Die Behebung.}

\begin{center}
{\small
\renewcommand{\arraystretch}{1.45}
\begin{tabular}{@{}P{3.0cm}P{2.2cm}P{2.6cm}@{}}
\toprule
\textbf{Weg} & \textbf{Dauer} & \\
\midrule
Einzeln & $\mathbf{5{,}1}$ s & \\
Alle Fallnummern auf einmal & $\mathbf{78}$ ms & $65\times$\\
Ein Aufruf für alles & $\mathbf{61}$ ms & \textcolor{gruen}{\textbf{fachlicher Zuschnitt}}\\
Nebenläufig einzeln & $340$ ms & \textcolor{mittel}{hilft, löst nicht}\\
\bottomrule
\end{tabular}}
\end{center}

\begin{center}
\kursiv{Die letzte Zeile ist die Versuchung: Nebenläufigkeit bringt Faktor fünfzehn und die Bündelung Faktor fünfundsechzig. Wer nebenläufig einzeln abfragt, hat das Problem verschleiert statt behoben --- und die Last auf der Gegenseite verzehnfacht.}
\end{center}

\begin{warnung}
\boxtitelr{Der Zuschnitt der Naht entscheidet}
\begin{center}
\renewcommand{\arraystretch}{1.5}
\begin{tabular}{@{}P{2.8cm}P{2.4cm}P{2.8cm}@{}}
\toprule
& \textbf{Technisch geschnitten} & \textbf{Fachlich geschnitten}\\
\midrule
Aufrufe je Ansicht & \textcolor{akzent}{$\mathbf{n{+}1}$} & \textcolor{gruen}{$\mathbf{1}$}\\
Wiederver\-wendbar & \textcolor{gruen}{ja} & \textcolor{mittel}{weniger}\\
Übertragene Menge & \textcolor{gruen}{genau} & \textcolor{mittel}{etwas mehr}\\
Passt zur Aufgabe & \textcolor{akzent}{\textbf{nein}} & \textcolor{gruen}{\textbf{ja}}\\
\bottomrule
\end{tabular}
\end{center}

\begin{center}
\begin{minipage}{0.88\textwidth}
\kursiv{Eine Naht, die \code{holeBefund(id)} anbietet, ist allgemein und erzwingt $n{+}1$. Eine, die \code{holeVerlauf(patient, groesse, zeitraum)} anbietet, ist enger und beantwortet die tatsächliche Frage in einem Aufruf --- Ousterhouts tiefe Schnittstelle aus Band 20, hier über das Netz.}
\end{minipage}
\end{center}

\begin{center}
\kursiv{Der Preis steht in der zweiten Zeile und ist geringer, als er wirkt: Eine Naht mit acht fachlichen Operationen deckt fast alles ab, was gebraucht wird. Wer hundert allgemeine Operationen anbietet, hat keine Schnittstelle, sondern eine Datenbank über das Netz.}
\end{center}
\end{warnung}
\end{merksatz}

\begin{aufgabe}
\begin{enumerate}[label=(\alph*)]
\item Verfolgen Sie eine langsame Anfrage.
\item Rechnen Sie die Wirkung der Bündelung aus.
\item Bestimmen Sie, wo man zuerst ansetzt.
\end{enumerate}

\loesung
\textbf{(a) Eine langsame Anfrage verfolgen.}

\textbf{Die Meldung.} Der Aufbau der Verlaufsansicht dauert bei manchen Patienten über fünf Sekunden.

\begin{center}
{\small
\renewcommand{\arraystretch}{1.45}
\begin{tabular}{@{}P{1.0cm}P{3.2cm}P{2.0cm}P{2.2cm}@{}}
\toprule
& \textbf{Abschnitt} & \textbf{Dauer} & \textbf{Anteil}\\
\midrule
$1$ & Oberfläche & $40$ ms & $1\,\%$\\
$2$ & Auswertung, Rechnung & $180$ ms & $4\,\%$\\
$3$ & \textbf{Aufrufe an Speicherung} & \textcolor{akzent}{$\mathbf{4{,}6}$ \textbf{s}} & \textcolor{akzent}{$\mathbf{90\,\%}$}\\
$4$ & davon Netz & $1{,}1$ s & \\
$5$ & davon Serialisierung & $\mathbf{0{,}9}$ s & \\
$6$ & davon Datenbank & $2{,}6$ s & \\
\bottomrule
\end{tabular}}
\end{center}

\begin{center}
\begin{minipage}{0.88\textwidth}
\kursiv{Neunzig Prozent in einem Abschnitt, und die Aufteilung darunter zeigt: Selbst wenn die Datenbank unendlich schnell wäre, blieben zwei Sekunden. Das ist der Befund, den ein Datenbankprofil allein nie liefert.}
\end{minipage}
\end{center}

\begin{center}
\kursiv{Die fünfte Zeile ist der überraschende Posten: Neunhundert Millisekunden nur für das Umwandeln in JSON und zurück --- bei dreihundertvierzig Einzelaufrufen mit je einem kleinen Ergebnis.}
\end{center}

\medskip
\textbf{(b) Die Wirkung der Bündelung.}

\begin{center}
{\small
\renewcommand{\arraystretch}{1.45}
\begin{tabular}{@{}P{2.4cm}P{1.8cm}P{1.8cm}P{2.2cm}@{}}
\toprule
\textbf{Maßnahme} & \textbf{Netz} & \textbf{Serial.} & \textbf{Datenbank}\\
\midrule
Vorher & $1{,}1$ s & $0{,}9$ s & $2{,}6$ s\\
Gebündelt & $\mathbf{4}$ ms & $\mathbf{31}$ ms & $\mathbf{43}$ ms\\
Faktor & $\mathbf{275}$ & $29$ & $\mathbf{60}$\\
\bottomrule
\end{tabular}}
\end{center}

\begin{center}
\begin{minipage}{0.88\textwidth}
\kursiv{Eine Änderung, drei Posten behoben. Der größte Faktor liegt beim Netz, weil dort die Verzögerung je Aufruf anfällt und nicht die übertragene Menge --- dreihundertvierzig mal drei Millisekunden gegen einmal vier.}
\end{minipage}
\end{center}

\begin{center}
\renewcommand{\arraystretch}{1.45}
\begin{tabular}{@{}P{2.8cm}P{2.2cm}P{2.4cm}P{1.8cm}@{}}
\toprule
& \textbf{Vorher} & \textbf{Nachher} & \\
\midrule
Gesamt & $\mathbf{5{,}1}$ s & $\mathbf{298}$ ms & $17\times$\\
Übertragene Menge & $1{,}2$ MB & $\mathbf{1{,}4}$ MB & \textcolor{mittel}{$+17\,\%$}\\
Zeilen geändert & --- & $\mathbf{24}$ & \\
Aufwand & --- & $\mathbf{4}$ h & \\
\bottomrule
\end{tabular}
\end{center}

\begin{center}
\kursiv{Vier Stunden Arbeit, vierundzwanzig geänderte Zeilen, Faktor siebzehn. Die zweite Zeile ist der Preis: siebzehn Prozent mehr Daten, weil der gebündelte Aufruf etwas mitliefert, was nicht jeder braucht.}
\end{center}

\medskip
\textbf{(c) Wo man zuerst ansetzt.}

\begin{center}
{\small
\renewcommand{\arraystretch}{1.45}
\begin{tabular}{@{}P{1.0cm}P{3.2cm}P{2.0cm}P{2.2cm}@{}}
\toprule
& \textbf{Schritt} & \textbf{Aufwand} & \textbf{Gewinn}\\
\midrule
$1$ & \textbf{Verfolgen, nicht raten} & \textcolor{gruen}{$2$ h} & --- \\
$2$ & \textbf{Aufrufe über Nähte bündeln} & \textcolor{gruen}{$1$ d} & \textcolor{gruen}{\textbf{$5$--$60\times$}}\\
$3$ & Datenbankindizes & $1$ d & $2$--$50\times$\\
$4$ & Zwischenspeicher & $2$ d & $2$--$20\times$\\
$5$ & Format wechseln & $3$ d & $2$--$5\times$\\
$6$ & Mehr Rechner & \textcolor{akzent}{laufend} & \textcolor{akzent}{$1$--$2\times$}\\
\bottomrule
\end{tabular}}
\end{center}

\begin{center}
\begin{minipage}{0.88\textwidth}
\kursiv{Die Reihenfolge ist nach dem Verhältnis von Gewinn zu Aufwand geordnet, und sie wird regelmäßig von unten begonnen. Schritt sechs ist der teuerste und der schwächste --- doppelte Rechnerzahl bringt bei einem $n{+}1$-Problem gar nichts, weil die Verzögerung nicht von der Rechenleistung abhängt.}
\end{minipage}
\end{center}

\begin{center}
\kursiv{Schritt zwei steht vorn, weil er in einem verteilten Werk fast immer trägt --- und weil er die einzige Maßnahme ist, die drei Kostenposten zugleich senkt: Netz, Serialisierung und Datenbank.}
\end{center}

\begin{warnung}
\boxtitelr{Die Verzögerung addiert sich, die Bandbreite nicht}
\begin{center}
\renewcommand{\arraystretch}{1.5}
\begin{tabular}{@{}P{2.8cm}P{2.4cm}P{2.8cm}@{}}
\toprule
& \textbf{Verzögerung} & \textbf{Bandbreite}\\
\midrule
Bei mehr Aufrufen & \textcolor{akzent}{\textbf{addiert sich}} & \textcolor{gruen}{unverändert}\\
Bei mehr Daten & \textcolor{gruen}{fast unverändert} & \textcolor{akzent}{steigt}\\
Durch schnellere Leitung & \textcolor{gruen}{kaum besser} & \textcolor{gruen}{\textbf{besser}}\\
Durch kürzere Strecke & \textcolor{gruen}{\textbf{besser}} & unverändert\\
\bottomrule
\end{tabular}
\end{center}

\begin{center}
\begin{minipage}{0.88\textwidth}
\kursiv{Die dritte Zeile überrascht regelmäßig: Eine zehnmal schnellere Leitung verkürzt die Verzögerung fast nicht, weil sie von der Laufzeit des Signals und der Zahl der Zwischenstationen bestimmt wird. Bandbreite kauft man, Verzögerung nicht.}
\end{minipage}
\end{center}

\begin{center}
\kursiv{Daraus folgt die praktische Regel: Wenige große Aufrufe schlagen viele kleine, auch wenn dabei mehr Daten übertragen werden. Siebzehn Prozent mehr Menge gegen Faktor siebzehn Zeit --- das ist der Handel aus Aufgabe (b).}
\end{center}
\end{warnung}
\end{aufgabe}
\vspace{-4pt}
\noindent Gregg\zit{11} zeigt, wie man Ab\-schnit\-te einer An\-fra\-ge trennt; Fowler\zit{8} sammelt die Bau\-for\-men, die $n{+}1$ ver\-mei\-den; Klepp\-mann\zit{15} ordnet die Ab\-wä\-gun\-gen zwischen Menge und Ver\-zö\-ge\-rung ein.
